% !PRINTED none
% Abschnitt-opening page (no running-head folio visible in any of the three bands):
% large upper blank area, then centred "Fünfter Abschnitt." / "Die Modulfunktionen."
% title block with a rule, followed directly by the "§ 51." heading and the
% start of its body text on the same leaf.
\begin{center}
\gesperrt{Fünfter Abschnitt.}

\textbf{Die Modulfunktionen.}

\rule{2cm}{0.4pt}

\textbf{\S\,51. \gesperrt{Die elliptischen Differentialgleichungen.}}
\end{center}

Die bisher definierten elliptischen Funktionen sind Funktionen der beiden Variablen $v$, $\omega$, und die Moduln $\varkappa$, $\varkappa'$, ferner $j(\omega)$, $K$, $K'$ sind von dem Periodenverhältnis $\omega$, das einen positiv imaginären Bestandteil haben muß, abhängig. Ebenso ist die Funktion $\wp(u)$ eine Funktion der drei Variablen $u$, $\omega_1$, $\omega_2$ und die Invarianten $g_2$, $g_3$ sind von $\omega_1$, $\omega_2$ abhängig. \gesperrt{Die Aufgabe der Umkehrung der elliptischen Integrale}, wie wir sie in \S\,14 formuliert haben, setzt aber voraus, daß in den elliptischen Funktionen $\varkappa^2$ (oder bei der \gesperrt{Weierstrass}schen Normalform $g_2, g_3$) beliebig gegebene Werte haben, und die vollständige Lösung dieser Aufgabe verlangt also, daß nicht $\omega$, sondern $\varkappa^2$ (bzw. $g_2, g_3$) als zweite unabhängige Variable auftritt, und daß $\omega$ durch diese ausgedrückt werde. Dieser Aufgabe werden die nächsten Betrachtungen gewidmet sein.

Wir gehen aus von der Aufgabe, ein System von Differentialgleichungen, welches wir das der \gesperrt{elliptischen Differentialgleichungen} nennen, zu integrieren:
\weq{1}{\begin{aligned}
\frac{dx}{dv} &= yz, \\
\frac{dy}{dv} &= -zx, \\
\frac{dz}{dv} &= -\varkappa^2xy,
\end{aligned}}
unter den Nebenbedingungen:
\weq{2}{\text{für } v=0 \text{ soll } x=0,\ y=1,\ z=1 \text{ sein.}}
